
@@M0@@; plaçons-nous par exemple sur @@M1@@. Alors @@M2@@ est défini par le module @@M3@@, et @@M4@@ par @@M5@@, le morphisme canonique envoyant @@M6@@ sur @@M7@@ pour tout @@M8@@. Alors @@M9@@ admet pour base les éléments @@M10@@ pour @@M11@@.

On a @@M12@@ ; posons @@M13@@.
Comme la suite des @@M14@@ est régulière dans @@M15@@, il en est de même de la suite @@M16@@ qui s'en déduit par une matrice inversible, et donc aussi de la suite @@M17@@ qui en est extraite (on est dans un cas noethérien). Par suite, @@M18@@ est un morphisme régulier.

Il faut donc montrer que les @@M19@@ engendrent l'Idéal de @@M20@@. Dans @@M21@@, nous noterons, pour @@M22@@, @@M23@@ l'élément @@M24@@ de @@M25@@, considéré comme étant de degré @@M26@@. Au-dessus de @@M27@@, @@M28@@, et @@M29@@ est défini par le morphisme qui envoie @@M30@@ sur @@M31@@. Dans ce morphisme, l'image de @@M32@@ est @@M33@@. Donc les @@M34@@ appartiennent à l'Idéal de @@M35@@.

Comme tout élément de @@M36@@ peut s'exprimer comme polynôme par rapport à @@M37@@, il suffit, pour montrer que les @@M38@@ engendrent l'Idéal de @@M39@@, de montrer qu'aucun élément non nul de @@M40@@ n'appartient à l'Idéal de @@M41@@, ce qui revient à montrer que l'homomorphisme
@@M42@@
est injectif, car @@M43@@ n'est pas diviseur de 0 dans @@M44@@. Si on gradue @@M45@@ en attribuant le degré @@M46@@ à @@M47@@, et le degré @@M48@@ aux @@M49@@, le noyau de l'homomorphisme considéré est un idéal gradué. Or il est clair que
